\chapter*{Kế hoạch thực hiện}
\addcontentsline{toc}{chapter}{Kế hoạch thực hiện}
\label{chap:kehoach}

Đồ án được thực hiện theo hai giai đoạn ứng với hai mốc nộp bài. \textbf{Giai đoạn 1} kéo dài ba tuần, kết thúc ở mốc nộp giữa kỳ ngày 15/10/2026, tập trung hoàn tất phần thiết kế còn dang dở và dựng nền hệ thống. \textbf{Giai đoạn 2} kéo dài mười một tuần, kết thúc ở mốc nộp cuối kỳ ngày 31/12/2026, dành cho hiện thực, kiểm thử và đánh giá.

Bảng~\ref{tab:kehoach} trình bày kế hoạch chi tiết. Cột \emph{Kết quả} ghi rõ mục nào trong báo cáo sẽ được hoàn thiện, để tiến độ luôn truy được về nội dung cụ thể thay vì chỉ là mô tả công việc.

\begingroup
\small
\renewcommand{\arraystretch}{1.25}
\setlength{\tabcolsep}{5pt}
\begin{longtable}{L{1.4cm} L{2.2cm} L{6.6cm} L{3.6cm}}
\caption{Kế hoạch thực hiện theo hai giai đoạn}
\label{tab:kehoach} \\
\toprule
\rowcolor{tblheadbg}
\textbf{Tuần} & \textbf{Thời gian} & \textbf{Công việc} & \textbf{Kết quả} \\
\midrule
\endfirsthead
\toprule
\rowcolor{tblheadbg}
\textbf{Tuần} & \textbf{Thời gian} & \textbf{Công việc} & \textbf{Kết quả} \\
\midrule
\endhead
\bottomrule
\endfoot
\multicolumn{4}{@{}l}{\cellcolor{tblheadbg}\textbf{Giai đoạn 1 --- tới mốc nộp giữa kỳ 15/10/2026}} \\
1 & 23--29/9 & Thiết kế API theo từng nhóm yêu cầu chức năng; thiết kế bộ test case bám các bảng ngưỡng & Mục 5.3.7, 5.3.9 \\
\rowcolor{tblaltbg}
2 & 30/9--6/10 & Use case tổng quát và theo tác nhân; sitemap; lược đồ tuần tự và lược đồ lớp & Mục 5.1.2, 5.3.2--5.3.4 \\
3 & 7--15/10 & Vẽ ba sơ đồ còn thiếu (ngữ cảnh, kiến trúc tổng thể, ERD); dựng khung dự án và môi trường phát triển; hiện thực module tài khoản, phân quyền và eKYC; hoàn thiện tóm tắt và lời cảm ơn & Mục 4.2.1, 5.3.1, 5.3.5, 6.1.1--6.1.3; phần mở đầu \\
\multicolumn{4}{@{}l}{\cellcolor{tblheadbg}\textbf{Giai đoạn 2 --- tới mốc nộp cuối kỳ 31/12/2026}} \\
\rowcolor{tblaltbg}
4--5 & 16--29/10 & Hiện thực module sản phẩm, lô hàng và nhật ký canh tác; sinh mã QR và trang truy xuất công khai; module quản lý chứng nhận & Mục 6.1.3 \\
6--7 & 30/10--12/11 & Hiện thực module thu mua và tồn kho theo lô; module bán lẻ, trong đó có cơ chế chống bán vượt tồn kho khi hai luồng giao dịch đồng thời & Mục 6.1.3 \\
\rowcolor{tblaltbg}
8--9 & 13--26/11 & Hiện thực module bán sỉ (báo giá, hợp đồng, giao theo đợt) và module thanh toán kèm cơ chế tạm giữ tiền & Mục 6.1.3 \\
10 & 27/11--3/12 & Hiện thực module vận chuyển và khiếu nại, đánh giá và uy tín, quản trị & Mục 6.1.3 \\
\rowcolor{tblaltbg}
11 & 4--10/12 & Hiện thực và huấn luyện hai mô-đun AI: gợi ý sản phẩm và dự báo giá & Mục 6.1.3, 6.1.4 \\
12 & 11--17/12 & Thực thi bộ test case; kiểm thử tải và kiểm thử tranh chấp tồn kho; rà soát phân quyền; sửa lỗi & Mục 6.2 \\
\rowcolor{tblaltbg}
13 & 18--24/12 & Khảo sát độ khả dụng với nhà cung cấp theo bộ công cụ tại Phụ lục B; đo hiệu năng, bảo mật và chất lượng mô-đun AI; điền kết quả đánh giá & Chương 7 \\
14 & 25--31/12 & Viết phần kết quả đạt được; rà soát toàn bộ báo cáo; chuẩn bị slide và bản demo & Chương 8; slide, demo \\
\end{longtable}
\endgroup

\noindent\textbf{Rủi ro tiến độ và phương án dự phòng.} Giai đoạn 2 phải hiện thực tám module nghiệp vụ và hai mô-đun AI trong mười một tuần --- đây là phần rủi ro nhất của kế hoạch. Hai biện pháp được chuẩn bị sẵn:

\begin{itemize}
    \item \textbf{Thứ tự hiện thực bám theo mức độ ưu tiên đã xác định} (Phụ lục M): tám nhóm yêu cầu ở mức Must được làm trước, ba nhóm Should và hai nhóm Could xếp sau. Nếu thiếu thời gian, phần cắt sẽ rơi vào nhóm Could --- cụ thể là tương tác người dùng và mô-đun dự báo giá --- chứ không ảnh hưởng tới các luồng nghiệp vụ cốt lõi.
    \item \textbf{Module khó nhất được xếp sớm.} Cơ chế chống bán vượt tồn kho khi một lô được giao dịch đồng thời ở hai luồng là ràng buộc kỹ thuật khó nhất (NFR1.2, NFR4.1) và được xếp vào tuần 6--7 thay vì cuối kỳ, để nếu phát sinh vướng mắc thì vẫn còn thời gian xử lý.
\end{itemize}

\noindent Tuần 14 không có phần dự phòng riêng, nên mọi chậm trễ ở giai đoạn 2 phải được xử lý ngay trong tuần phát sinh bằng cách cắt hạng mục ưu tiên thấp, không dồn về cuối.
